\section{Skins and grounded extraction}
\label{app:skins}
\paragraph{Skins.} Each skin names the parties, the artifact, four issues and three options per issue; the generator still draws which option scores what, and under a skin the letters are kept (``A (an in-process cache)'').
\begin{description}\small
  \item[software design] Developer 1 and 2 agree on the design of a shared module; ``you each own one feature of the same service, and both features go through this module''. Storage layer: an in-process cache, the existing SQL tables, a separate key-value store. Call style: synchronous calls, callbacks, an event queue. Error handling: raise exceptions, return result objects, log and continue. Configuration: constants in code, a config file, environment variables.
  \item[team plan] Team lead 1 and 2 agree on a quarterly plan for both teams; ``your two teams share one budget, one office and one review process''. Hiring: two senior engineers, one senior and two junior engineers, four junior engineers. Office days: Monday to Wednesday, Tuesday to Thursday, Wednesday to Friday. Tooling budget: cloud credits, workstations, training. Review cadence: weekly, every two weeks, monthly.
  \item[release plan] Product owner 1 and 2 agree on a joint release plan; ``your two products ship together in one release''. Release date: end of March, end of April, end of May. First platform: web, iOS, Android. Pricing: a free tier plus a paid plan, paid plans only, usage-based pricing. Launch region: North America, Europe, Asia.
\end{description}

\paragraph{Grounded extraction.} For each of ParallACT's 85 contradictory pairs the extractor (\texttt{gpt-6-luna}, low reasoning) receives both tickets, the explanation of why their changes collide and the reconciliation, and returns JSON with two to four decisions; for each decision a name, three options labeled as the one ticket A needs, the one ticket B needs and a middle ground, and each side's importance from 1 to 5; a role phrase per developer; and a neutral context paragraph of two to three sentences. Validation requires the middle option to differ from both needs, distinct option names, and at least one decision on which the two needs differ; a failed extraction is retried once with the validation errors. A first version that asked for the options in a fixed order was misordered in about half of the decisions by the model with reasoning off; labeling the three roles explicitly extracted all 85 pairs (4 on the second attempt) for \$0.07. A decision on which the two needs coincide is a shared issue; importances are normalized to 100 per side; option letters are shuffled per instance with a fixed seed. \tbd{Two example extractions, and the result of the human spot-check.}

